\documentclass[tikz,border=3mm,multi=tikzpicture]{standalone}
\usepackage{eleve-matematica-ef1}

\begin{document}

% FIGURA 1 — fig-01-soma-mesmo-denominador
\begin{tikzpicture}[matematica figura, x=1cm, y=1cm]
  \path[use as bounding box] (-0.15,-0.20) rectangle (8.15,3.45);
  \foreach \i in {0,...,3}
    \draw[matematica contorno] ({0.10+2*\i},1.10) rectangle ({2.10+2*\i},3.10);
  \fill[matemalaranjaclaro] (0.14,1.14) rectangle (2.06,3.06);
  \fill[matemazulclaro] (2.14,1.14) rectangle (6.06,3.06);
  \draw[matematica contorno] (0.10,1.10) rectangle (8.10,3.10);
  \foreach \x in {2.10,4.10,6.10}
    \draw[matematica divisao] (\x,1.10)--(\x,3.10);
  \node[matematica rotulo, text=matemalaranja] at (1.10,0.52) {almoço: $\frac14$};
  \node[matematica rotulo, text=matemazul] at (4.10,0.52) {jantar: $\frac24$};
  \node[matematica resultado] at (7.10,0.52) {total: $\frac34$};
\end{tikzpicture}

% FIGURA 2 — fig-02-subtracao-mesmo-denominador
\begin{tikzpicture}[matematica figura, x=1cm, y=1cm]
  \path[use as bounding box] (-0.10,-0.15) rectangle (8.80,6.10);
  \foreach \x/\n in {0.60/antes,6.10/depois}{
    \draw[matematica contorno, rounded corners=5pt] (\x,0.55) rectangle ({\x+2.10},5.25);
    \foreach \j in {1,...,4}
      \draw[matematica divisao] (\x,{0.55+0.94*\j})--({\x+2.10},{0.55+0.94*\j});
    \node[matematica rotulo] at ({\x+1.05},5.72) {\n};
  }
  \fill[matemazulclaro] (0.64,0.59) rectangle (2.66,4.27);
  \fill[matemazulclaro] (6.14,0.59) rectangle (8.16,3.33);
  \draw[matematica contorno, rounded corners=5pt] (0.60,0.55) rectangle (2.70,5.25);
  \draw[matematica contorno, rounded corners=5pt] (6.10,0.55) rectangle (8.20,5.25);
  \foreach \x in {0.60,6.10}
    \foreach \j in {1,...,4}
      \draw[matematica divisao] (\x,{0.55+0.94*\j})--({\x+2.10},{0.55+0.94*\j});
  \draw[-{Stealth[length=3.2mm]},matematica destaque] (3.20,2.90)--(5.55,2.90);
  \node[matematica resultado] at (4.38,3.45) {$-\frac15$};
  \node[matematica valor] at (1.65,0.05) {$\frac45$};
  \node[matematica valor] at (7.15,0.05) {$\frac35$};
\end{tikzpicture}

% FIGURA 3 — fig-03-completar-inteiro
\begin{tikzpicture}[matematica figura, x=1cm, y=1cm]
  \path[use as bounding box] (-0.15,-0.10) rectangle (8.35,4.00);
  \foreach \i in {0,...,3}
    \draw[matematica contorno] ({0.10+1.60*\i},1.00) rectangle ({1.70+1.60*\i},2.60);
  \foreach \i in {0,1,2}
    \fill[matemalaranjaclaro] ({0.14+1.60*\i},1.04) rectangle ({1.66+1.60*\i},2.56);
  \fill[matemazulclaro] (4.94,1.04) rectangle (6.46,2.56);
  \draw[matematica contorno] (0.10,1.00) rectangle (6.50,2.60);
  \foreach \x in {1.70,3.30,4.90}
    \draw[matematica divisao] (\x,1.00)--(\x,2.60);
  \node[matematica rotulo, text=matemalaranja] at (2.50,3.25) {parte marcada: $\frac34$};
  \node[matematica rotulo, text=matemazul] at (5.70,3.25) {falta: $\frac14$};
  \draw[-{Stealth[length=3mm]},matematica destaque] (6.80,1.80)--(7.65,1.80);
  \node[matematica resultado] at (7.85,1.80) {$1$};
\end{tikzpicture}

\end{document}
